% Pursuing Stacks (cotes 134-2 to 134-8): archive cote 134-8, site page 30 (PDF page 31), typescript page 584.
% typescript. Draft compared: ps7.tex:1392-1455.
% AI-transcribed by Claude Opus 5.5 (claude-opus-5-5, adaptive thinking, effort high), Message Batches API, batch ps4 (msgbatch_012Aoc1NuskATPsyBeKSYK7J), 2026-11-12; not yet checked by eye.
\scanpage{134-8/30}
\mspage{584}

where now ${}_{b\backslash}A$  is the category of pairs
\[ (a,u) \ \text{ with } \ u: b \to f_{o}(a) \quad . \]
We'll have to assume \unsure{npw|now} that both sheaves of rings $\underline{O}_A$ , $\underline{O}_B$  are con-%
stant,  and correspond to the same ring k.  \xout{Thexcommutationxofxf} Thus
Modules on A and B are just contravariant functors from these catego-%
ries to the category  $Ab_k$  of \xout{all} k-modules,  and accordingly $f_!$
\inserted[typed]{left adjoint \xout{to}}(\xout{dual} to the composition functor $f^*$ )  may be computed by a well known
formula,  involving direct limits on the categories ${}_{b\backslash}A$ :
\[ (7) \qquad f_!(F)(b) \simeq \varinjlim_{b\backslash A} F(a) \qquad , \]
where the limit in the second member is relative to the \xout{funct} composi-%
tion
\[ {}_{b\backslash}A \longrightarrow A \xrightarrow{\ F\ } Ab_k \quad . \]
(NB This formula is dual to the formula for $f_*$ ,  the right adjoint of
$f^*$ ,  closer to intuition  - to mine at any rate -  because the $\varprojlim$ in
\[ f_*(F)(b) \simeq \varprojlim_{A_{/b}} F(a) \]
may be ``visualized'' as \xout{\unsure{Asectionsx}}  the \xout{setxof} set of ``sections'' of
F over $A_{/b}$ \unsure{.}) From this formula (7) follows at once commutation of
$f_!$  with colocalization\unsure{.}  To get the corresponding result for $Lf_!$ ,
we have only to use the fact that the relevant inverse image functor
(corresponding to  ${}_{b\backslash}A \to A$ ) transforms projective\deleted{\illegible{1}} modules into
% NOTE: from here to the foot of the page the text is typed on a strip pasted over the page; the older typed text underneath shows through faintly and is not transcribed
projective modules.  As the map  ${}_{b\backslash}B \longrightarrow B$  is a \uline{cofibering functor}
\uline{with discrete fibers},  so is the map
\[ h: A' = {}_{b\backslash}A \longrightarrow A \]
deduced by base change.  Now for any \inserted{such} functor  $h: A' \to A$ between small
categories,  the inverse image functor
\[ h^* : A^{\wedge} \longrightarrow A'^{\wedge} \ , \quad \text{or} \quad h^*_k : A^{\wedge}_k \to A'^{\wedge}_k \]
carries \inserted{indeed} projectives into projectives.  This statement is dual formally
to the corresponding statement for injectives,  valid when we make on
h the dual assumption of being \uline{fibering with discrete fibers}  -  in the
latter case the \hand{(}well known\hand{)} proof comes out formally from the fact
that the left adjoint functor $h_!$ or $h^k_!$ carries monomorphisms into
monomorphisms \deleted{\unsure{i.e. is exact}}  -  a fact which we used in section 135
in the form  $Lh^k_! = h^k_!$ (in case $k=\underline{Z}$).  In the present case the proof is
essentially the dual one  -  as a matter of fact, as there are enough

\margin{30}
